\documentclass[11pt,a4paper]{article}

% ---------------------------------------------------------------------------
% Preamble
% ---------------------------------------------------------------------------
\usepackage[utf8]{inputenc}
\usepackage[T1]{fontenc}
\usepackage{amsmath}
\usepackage{amssymb}
\usepackage{graphicx}
\usepackage{booktabs}
\usepackage{array}
\usepackage{longtable}
\usepackage[margin=1in]{geometry}
\usepackage{hyperref}

\hypersetup{
  colorlinks=true,
  linkcolor=black,
  citecolor=black,
  urlcolor=blue,
  pdftitle={Locating and Auditing a CompTIA Security+ Study Module in a Personal Obsidian Vault: A Retrieval and Quality-Assurance Case Report},
  pdfauthor={Author Placeholder}
}

\setlength{\parskip}{0.4em}
\setlength{\itemsep}{2pt}
\setlength{\parsep}{0.2em}
\renewcommand{\arraystretch}{1.15}

% ---------------------------------------------------------------------------
% Document
% ---------------------------------------------------------------------------
\title{Locating and Auditing a CompTIA Security+ Study Module\\
in a Personal Obsidian Vault:\\
A Retrieval and Quality-Assurance Case Report}
\author{Author Placeholder\\
\small Affiliation Placeholder}
\date{25 September 2026}

\begin{document}

\maketitle

\begin{abstract}
\noindent
Personal knowledge vaults are frequently used as private study corpora, yet the
provenance and internal consistency of the material stored in them are rarely
audited. This paper reports a retrieval-and-audit case study conducted on a
single personal Obsidian vault. A case-insensitive recursive search across the
vault root returned exactly one matching study artifact: a Security+ textbook by
Mark Ciampa, held in two distinct forms---a source PDF and a directory of
fifteen extracted per-module notes. The study scope was then narrowed to
Module~7, ``Identity and Access Management,'' associated in the vault with the
SY0-701 objective domain. That module was read in full at 471 lines and is
organised around three stated objectives: authentication credential types,
authentication best practices, and access control schemes. An initial quality
audit surfaced two defects: an internal contradiction between a stated industry
password-expiry practice and an explicit recommendation against expiry, and a
malformed review question offering two textually identical answer options. The
study then adopted a drill protocol in which twenty-five examination items drawn
from the module were each scored and accompanied by a line-level citation into
the source note. All twenty-five items were resolved against the note. This
report presents that item analysis descriptively, enumerates the rejected
distractor options alongside the accepted selections, and preserves three
defects as open findings rather than resolving them silently. No pass rates,
empirical effect sizes, or external literature are claimed; the evidence base
is one module of one note in one vault.
\end{abstract}

% TODO: the drill protocol's line-level citation apparatus is part of the
% working record but its locators are not reproduced here, because the
% distilled summary did not carry them. They must be re-attached from the
% working record before any version of this paper is cited.

\section{Introduction}
\label{sec:intro}

\subsection{Motivation}

A personal knowledge vault accumulated over years is a useful study corpus and
a poor publication. Notes in such a vault are typically transcribed or
machine-extracted from copyrighted source material, summarised under time
pressure, and revisited months later. Three failure modes follow from that
lifecycle. First, \emph{retrieval ambiguity}: a reader who knows a study title
may not know whether the vault holds it as a source document, as extracted
notes, or as both, and a naive filename search can silently return a partial
match. Second, \emph{latent contradiction}: summarisation can preserve both a
statement of prevailing industry practice and its own contrary recommendation
without the summariser noticing, leaving a note that cannot be answered
consistently. Third, \emph{malformed assessment}: transcription can produce
review questions whose options are duplicated, rendering the item unanswerable
regardless of the reader's knowledge.

These are quality-assurance problems, not pedagogy problems. They are usually
undetected because nothing in the workflow ever asks whether a note is
internally coherent or whether its review questions are well-formed.

\subsection{Research questions}

The study was organised around four questions, each answerable from the vault
alone:

\begin{enumerate}
  \item \textbf{RQ1 (Retrieval).} What CompTIA study material is present in the
        vault, and in what forms?
  \item \textbf{RQ2 (Structure).} How is the retrieved material organised, and
        which unit is suitable for close audit?
  \item \textbf{RQ3 (Defect detection).} What internal defects are present in
        the audited unit?
  \item \textbf{RQ4 (Grounding).} Can a set of twenty-five examination items
        drawn from the audited unit be resolved unambiguously against the
        note's own text, and which items cannot be?
\end{enumerate}

\subsection{Contribution}

The contribution is a worked, reproducible account of a retrieval-then-audit
procedure applied to one real study module, together with an item-level
resolution record that makes the note's coverage inspectable and preserves the
defects found. Section~\ref{sec:methods} gives the procedure,
Section~\ref{sec:results} the findings, and Section~\ref{sec:discussion} the
limits on what this evidence can support.

\section{Background and Material}
\label{sec:background}

\subsection{The retrieval event}

The study originated from an operational request: locate a CompTIA study book
within the personal Obsidian vault. A case-insensitive recursive search of the
vault root was issued against the target title.

The search returned exactly one matching title: a Security+ textbook by Mark
Ciampa. The matched material exists in the vault in two distinct forms.

\begin{description}
  \item[Source PDF.] The book in its original form as a single document.
  \item[Extracted notes.] A directory of fifteen per-module notes derived from
        that book, one directory entry per module.
\end{description}

% TODO: the vault path, the PDF filename, the PDF byte size, and the
% bibliographic details of the Ciampa volume (edition, publisher, year, ISBN)
% exist in the working record. They are omitted here because the distilled
% summary did not carry them; a citation-complete version must add them from
% the vault itself rather than from memory.

The dual representation is the operationally important finding of RQ1. A reader
who searched only for the PDF would have found one artefact and concluded the
book was available only in source form; a reader who searched only for note
directories would have found fifteen entries and concluded the opposite. The
book is in the vault in both forms simultaneously, and the two forms carry
different audit exposure: the source PDF is fixed, whereas the fifteen notes
are individually edited and individually fallible.

\subsection{Audited unit}

The user narrowed scope to Module~7, ``Identity and Access Management,''
recorded in the vault against SY0-701. This unit was read in full at 471 lines.
It is organised around three objectives:

\begin{enumerate}
  \item authentication credential types;
  \item authentication best practices;
  \item access control schemes.
\end{enumerate}

The module is therefore a natural audit target: three separable objective
areas, a size small enough to read end to end rather than sample, and a domain
(identity and access) whose canonical vocabulary is precise enough that a
distractor option can be adjudicated against the note's own wording rather
than against opinion.

\section{Methods}
\label{sec:methods}

\subsection{Design}

The study is a single-case, document-level audit with an embedded
verification drill. It is descriptive and non-statistical. There is no control
group, no comparison corpus, no participant sample, and therefore no basis for
an effect size, a pass rate, or a significance claim; none is reported.

\subsection{Procedure}

The procedure had four stages, executed in this order.

\textbf{Stage 1 --- Retrieval.} A case-insensitive recursive search was issued
across the vault root. The case-insensitivity was deliberate: the user
requested the book by informal description rather than by exact filename, and
the search had to be robust to title-casing differences between the source
filename and the note directory name.

\textbf{Stage 2 --- Scope selection.} Both retrieved forms were inventoried and
the fifteen-module note directory enumerated. The user then selected a single
unit for close audit, which fixed Module~7 as the unit of analysis and the
three objectives above as its top-level partition.

\textbf{Stage 3 --- Defect audit.} The selected module was read in full, not
sampled. Reading end to end is the only way to detect a contradiction between
two passages that are far apart in the file; sampling would have had a
substantial probability of encountering only one side of it. Two defects were
recorded at this stage (Section~\ref{sec:results:defects}).

\textbf{Stage 4 --- Verification drill.} The user supplied twenty-five
examination items drawn from the module. A drill protocol was adopted under
which every item is adjudicated by two outputs together: a \emph{score}
(accepted selections and rejected selections) and a \emph{citation} of the
supporting lines in the source note. The citation is what makes the drill an
audit rather than an opinion: any score can be challenged by pointing at the
note, and any passage of the note can be turned back into an item.

\subsection{Scoring conventions}

Each item's resolution is recorded as an ordered triple:
\begin{enumerate}
  \item the \textbf{accepted} option or options, i.e.\ those the note supports;
  \item the \textbf{rejected} option or options, i.e.\ those the note does not
        support, together with the reason; and
  \item the \textbf{status}: \textsc{resolved} when the note determines the
        answer, or \textsc{ambiguous} when the item demands a selection the
        note does not determine.
\end{enumerate}

A distractor was recorded as rejected only when the note gave grounds for
rejection. Where rejection turned on a single stated property --- for example,
that peppering leaves the hashing function untouched, or that an access control
list's documented limitation is inefficiency --- the property is carried in the
reject record rather than asserted bare, so that the rejection is auditable
against the same note.

\subsection{Taxonomy used for reporting}

% TODO: the mapping of the twenty-five items onto the three module objectives
% is the analyst's grouping, made for reporting convenience. The distilled
% summary does not state the item-to-objective mapping. It is offered as a
% reporting device and should be read as such, not as a source fact.
The three module objectives are used as a reporting taxonomy for the
twenty-five items. The grouping is the analyst's, introduced to make the
coverage of the drill inspectable; it is not a partition stated by the note.

\section{Results}
\label{sec:results}

\subsection{Retrieval outcome (RQ1)}

The case-insensitive recursive search returned a single matching title: a
Security+ textbook by Mark Ciampa, present both as a source PDF and as a
directory of fifteen extracted per-module notes. No second CompTIA title was
matched. The corpus under audit is therefore a single book's two
representations, not a multi-title collection.

\subsection{Unit structure (RQ2)}

The audited unit is Module~7, ``Identity and Access Management'' (SY0-701),
read in full at 471 lines, partitioned across the three objectives stated in
Section~\ref{sec:background}. Fifteen modules exist in the note directory;
one was audited.

\subsection{Defects found on the initial audit (RQ3)}
\label{sec:results:defects}

The full read surfaced two defects. Both are properties of the note, not of the
reader.

\textbf{D1 --- Password-expiry contradiction.} The note states an industry
password-expiry practice and, elsewhere, explicitly recommends against
expiry. The two statements cannot both be operative guidance. This is a
substantive contradiction about what to do with passwords, and it is the more
consequential of the two defects, because a reader who follows the wrong half
of it implements a control the other half argues against. The defect is
recorded here as an open finding; the drill did not settle it, and the summary
does not indicate which statement the note intends to govern.

\textbf{D2 --- Malformed review question.} The note contains a review question
that presents two textually identical answer options. An item whose options are
duplicates is unanswerable: a correct reader and an incorrect reader select the
same string. This is a mechanical defect of transcription, and it is recorded
as an open finding for the same reason as D1.

Both defects are of a kind that sampling would plausibly have missed, which is
the empirical justification for the full read in Stage~3.

\subsection{Resolution of the twenty-five drill items (RQ4)}
\label{sec:results:drill}

All twenty-five items were resolved against the note. The full item register,
with accepted and rejected options for each item, is given as
Appendix~\ref{app:register}. The substantive findings are grouped below by the
mechanism or scheme each item exercised.

\subsubsection{Digest protection: salting, stretching, and masking}

Salting was identified as the mechanism that guarantees distinct digests for
identical plaintext. The grounding is the per-user salt: a per-user salt
guarantees divergence between users whose plaintexts coincide, whereas
peppering does not, because peppering leaves the hashing function untouched and
therefore leaves the digest unaltered. On this reasoning both symmetric and
asymmetric encryption were rejected as answers to the same item: neither alters
the digest, so neither can be the mechanism that makes identical plaintexts
hash differently. This is the sharpest single item in the drill, because the
rejections are forced by one stated property rather than by preference.

Key stretching and Argon2 were selected jointly as the mechanism protecting a
digest, the two together answering a two-part item on digest protection.

The complementary property was tested in the other direction. Where a password
is built by a statistical rule rather than by memorisation --- a mask
constructed from statistical properties of the target's own choices --- the
resulting attack is a rule attack, not a generic offline attack. Password
attacks were therefore partitioned into rule construction and offline
enumeration, and the rule-attack item was scored as the rule attack.

A separate item examined a sample password and asked what its structure
reveals. Predictable substitution-and-append structure was recognised, and the
password scored as weak on that basis: a structure that is predictable is a
structure an attacker's generator already produces.

\subsubsection{Attack taxonomy: yield, platform, and escalation}

The drill scored a technique as \emph{high-outcome} or \emph{low-outcome}
against the note's own classification, and that classification was enforced
rather than second-guessed. Two consequences follow. High-outcome cracking was
correctly tied to stolen digest files and to reverse engineering --- that is,
to obtaining the stored digest and then attacking it offline. Online brute force
was excluded on the grounds that the note classes it as low-outcome, which also
motivates account-lockout avoidance as a target technique for an adversary. The
mirror case is password spraying, whose selected motivations were account-lockout
avoidance and alarm avoidance, and whose low-outcome classification was likewise
enforced. Spraying's low yield is not a contradiction of its appeal; it is
precisely what makes it attractive, since a technique that evades lockout and
alarms while producing few successes is cheap to run and quiet to run.

The taxonomy also had to be checked for logical containment. Offline brute force
and dictionary attack are related by subsumption: a dictionary attack is a
restricted case of offline brute force, so an item asking how the two are
related is answered by that containment, and not by the inverted claim. The
inverted online-offline claim and a malformed comparison option were both
rejected.

Yield was additionally ranked, not merely classified. The note documents an
eight-stage escalation sequence for password attacks, and dictionary attack was
scored as the highest-yield early tool in that sequence, ahead of hybrid,
mask, and terminal brute force.
% TODO: the source note contains the full eight-stage sequence. Only the
% stages named in the distilled summary are reported here; the remaining
% stages must be transcribed from the note before the sequence is described in
% full.

Finally, the attacker's hardware was scored: parallel graphics processing was
identified as the desirable platform for brute force, on the reasoning that a
brute-force attack is a massively parallel workload and the platform that
parallelises it best is the platform worth wanting.

\subsubsection{Credential factors: possession, knowledge, and behaviour}

The possession element of authentication was the single most productive source
of item adjudications. A gate officer who admits a person on recognition was
scored as a failure to enforce the possession element: recognition establishes
something about the person, and possession is about the thing held. A related
distractor was rejected for the same reason --- namely, the claim that a
particular control satisfies the possession element when the note's reasoning
establishes it as a knowledge mechanism.

Knowledge-based possession was identified as the requirement for cognitive
biometrics, and the cognition difficulty was the point: a cognitive biometric is
one whose imitation is hard, so the control works because possession of the
answer is not possession of the answer \emph{alone}. The relative difficulty of
imitating cognitive biometrics was selected, and an absolute spoof-resistance
claim was rejected, because the note supports a comparison and not a guarantee.
Voice recognition was selected on the same item. The categorical assignment of
the picture-password scheme was rejected, and a spoof-resistance claim attached
to it was rejected with it.

The contrast the note draws is between the biological and the outwardly
exhibited, and genetic scanning was mapped accordingly: genetic scanning was
scored as the \emph{biological} element, not the outwardly exhibited one, since
a genome is not a behaviour that a subject performs in front of a sensor.

Keystroke dynamics were credited with requiring neither specialized hardware
nor additional user steps, which is what distinguishes a behavioural biometric
from the other two families in this item set.

Hardware factors were the remaining ground. Hardware password keys were
recommended against, on the grounds of malware exposure: a physical token's
advantage is that it cannot be typed into a hostile form, and malware is
precisely the threat that would have it typed. A hardware password manager
protected by a master password and a secret key file was classified as a
\emph{user} password manager rather than enterprise vaulting, the distinction
being between a product that stores one person's credentials and an
organisational facility that holds many principals' secrets.

In the federated direction, asymmetric cryptography was identified as the
signing basis for assertion-based federation, since an assertion is only
meaningful to a relying party if it is signed by something the relying party can
verify. In the phishing-resistant direction, passkeys were credited with
hardware storage and multi-factor combination, and security keys were preferred
over short-message authentication on the single ground of the absence of
transmitted codes --- there is no code to intercept, relay, or social-engineer.
The one-time-password item distinguished time-based from event-based
generation and is discussed as an open finding below.

\subsubsection{Password policy as practice}

Three policy elements were scored as sound practice: password length, reuse
prohibition, and minimum age. Two were rejected: expiration and an
eight-character complexity floor. Reuse prohibition appears a second time in
the drill, where it was scored as the sole defence against credential
stuffing --- the sole defence, because credential stuffing is defined by reuse
across breached accounts, so any policy permitting reuse leaves the
characteristic intact.

The two objectives were also reconciled conceptually. Identity proofing and
authentication were identified as the dual referents of the ``identity''
component of identity and access management: proofing is establishing that the
claimant is who they say they are, and authentication is the subsequent
control acting on that established identity.

\subsubsection{Access control schemes}

Rule-based role assignment was credited with dynamic role assignment, the
defining property of a role being that entitlements attach to the role and
therefore follow the user as assignments change. Attribute-based control was
credited with flexibility in combining attributes, the corresponding property
being that a decision is a function of several attributes at once rather than of
one role name. A user-modifiability claim and a conditional-structure claim were
rejected against the two schemes respectively.

Access control lists were matched to inefficient per-access permission
checking, the grounding being the note's own documented characterisation of
access control list limitations as inefficient and enterprise-ungainly. The
decision criterion is instructive: the item is not about whether access control
lists are correct, but about what their weakness is, and the note's stated
weakness is efficiency at scale.

\subsection{Coverage of the drill}

\begin{table}[h]
\centering
\caption{Items by module objective, using the reporting taxonomy of
Section~\ref{sec:methods}. The grouping is the analyst's.}
\label{tab:coverage}
\begin{tabular}{llr}
\toprule
\textbf{Objective} & \textbf{Representative item types} & \textbf{Items} \\
\midrule
Credential types & possession, knowledge, behaviour, storage, federation,
  passkeys, biometrics & 12 \\
Best practices & digest protection, password policy, attack taxonomy,
  password structure & 11 \\
Access control schemes & role-based assignment, attribute-based control,
  access control lists & 2 \\
\midrule
\textbf{Total} & & \textbf{25} \\
\bottomrule
\end{tabular}
\end{table}

Table~\ref{tab:coverage} shows the drill is weighted toward credential types
and authentication best practices, with access control schemes represented by
two items. Two items in the third group is a thin basis for any claim about
how thoroughly the access control scheme material was assessed, and the
imbalance should be read as a property of the item set the user supplied, not
as a property of the module.

\subsection{Open findings}
\label{sec:results:open}

Three items of unfinished business remain and are preserved rather than
resolved.

\begin{description}
  \item[O1 (D1).] The password-expiry contradiction in the module is
        unresolved. The note states an industry expiry practice and separately
        recommends against expiry, and the drill supplied no basis for deciding
        which governs.
  \item[O2 (D2).] The malformed review question with two textually identical
        options is unrepaired. Until one option is changed or the item is
        removed, it cannot discriminate between a correct and an incorrect
        reader.
  \item[O3.] The one-time-password item demanded two selections, but the note
        supports only one. The item as written is not answerable as a
        two-answer item, and the drill recorded the supported selection without
        manufacturing a second.
\end{description}

O1 and O2 are defects of the note and belong to the note's author. O3 is a
defect of the item as drafted rather than of the note, and is arguably the most
informative of the three, because it shows a distinct failure mode: not that
the source is silent, but that the question's expected answer count exceeds
what the source can support.

\section{Discussion}
\label{sec:discussion}

\subsection{What the case supports}

The case supports four claims, each of which follows from the recorded
procedure rather than from interpretation.

\textbf{Retrieval requires dual-form accounting.} The study material existed
simultaneously as a source PDF and as fifteen per-module notes. Any inventory
that reports only one form gives a false picture of what is available, and
gives it in whichever direction the search happened to look.

\textbf{Defects that matter are found by full reads.} Both defects found in
this module --- the password-expiry contradiction and the duplicated-option
review question --- are of types that a sampled read would plausibly have
missed, since each is a relation between two locations rather than a property
of either location. This is a small observational basis, one module, two
defects, and it supports only the weak form of the claim: no defect found here
was visible locally.

\textbf{Note-grounded drilling is adjudicable.} All twenty-five items were
resolved against the note, and a substantial share of the resolutions were
forced by a single stated property rather than by judgement: peppering leaves
the hashing function untouched, so it cannot alter a digest; symmetric and
asymmetric encryption do not alter digests either; a dictionary attack is a
restricted case of offline brute force; genetic scanning is biological and not
outwardly exhibited. Where the note is precise, item adjudication is
deterministic. The value of the citation requirement is visible in exactly these
items, because the reason for the score is the note line, not the analyst.

\textbf{Adversary economics unify the attack items.} The apparently scattered
attack items --- stolen digest files, reverse engineering, online brute force,
password spraying, lockout avoidance, alarm avoidance, dictionary-first
escalation, GPU platform --- resolve into a single account of attacker
economics: the note's outcome classification is the objective, and the items
that scored correctly were the ones that respected it. The spray and lockout
items in particular are only explicable once outcome class and cost-to-run are
read as the governing variables.

\subsection{What the case does not support}

The case does not support claims about learning outcomes, retention, exam
readiness, or the adequacy of the note as study material. No reader was
measured; the twenty-five items were adjudicated by an analyst against text, not
answered by a learner under time pressure. The count of twenty-five resolved
items is a measure of the note's adjudicability, not of the note's teaching
quality, and it must not be read as a pass rate.

The case is a single case. One module, one note, one analyst, one vault. The
taxonomy of Table~\ref{tab:coverage} is an analyst's construct, and the
two-item representation of access control schemes is too thin to characterise
that objective area. No comparison was made against the source PDF's own
coverage of the same material, so nothing here shows that the notes are
complete, only that they were sufficient to resolve the items supplied.

The \emph{low-outcome} and \emph{high-outcome} labels are the note's
classification, adopted as the scoring standard. Where the drill rejected an
option for being low-outcome --- online brute force, password spraying --- the
justification is fidelity to the source, not independent endorsement of the
source's risk assessment. An analyst who rejects those options for a different
reason would obtain the same score with a different rationale, and the citation
requirement is what keeps the two apart.

\subsection{Practical implications}

Three implications follow for anyone maintaining study material in a vault.
Contradictions should be repaired at the point of detection rather than deferred,
because a note that contradicts itself will be read with equal confidence in
both directions. Review items should be checked for well-formedness
mechanically, since duplicate options are trivially detectable and silently
destroy an item's discriminating power. And where a note is used to generate
examination items, the number of selections an item demands should be checked
against what the note actually determines, because an item demanding two
selections from a source that supports one is a defect of the item and will be
misattributed to the reader who cannot answer it.

\section{Conclusion}

A case-insensitive recursive search of a personal Obsidian vault located one
CompTIA study title, a Security+ textbook by Mark Ciampa, present both as a
source PDF and as a directory of fifteen extracted per-module notes. The
narrowed unit of study, Module~7 ``Identity and Access Management'' (SY0-701),
was read in full at 471 lines across three objectives. That read found two
defects: a contradiction between a stated industry password-expiry practice and
an explicit recommendation against expiry, and a review question offering two
textually identical options. Under a drill protocol in which each answer is
scored together with a line-level citation into the source note, all twenty-five
examination items drawn from the module were resolved against the note, covering
credential types, digest protection, password policy, the attack taxonomy, and
access control schemes, with passkeys, one-time passwords, asymmetric signing,
and cognitive biometrics each adjudicated from stated properties rather than
from preference. Three matters remain open: the expiry contradiction, the
malformed review item, and a one-time-password item whose demanded answer count
exceeded what the source supports. The case shows that a private vault is
auditable material: its retrieval is reproducible, its contradictions are
findable, and its review items are checkable, provided the audit reads the
module end to end and refuses to resolve what the source leaves undetermined.

\section*{References}
\addcontentsline{toc}{section}{References}

\begingroup
\renewcommand{\labelenumi}{[\arabic{enumi}]}
\begin{enumerate}
  \item Ciampa, M. \emph{Security+} (textbook by Mark Ciampa), as held in the
        author's personal Obsidian vault in two forms: a source PDF and a
        directory of fifteen extracted per-module notes. This is the sole
        substantive source for the present study.
        % TODO: full bibliographic record (exact title, edition, publisher,
        % year, ISBN) must be transcribed from the vault artefact or the
        % source PDF title page. It is deliberately not supplied from memory.
  \item \texttt{Module 7: Identity and Access Management} (SY0-701), extracted
        per-module note, 471 lines, personal Obsidian vault. Unit of analysis
        for the present study; the source of all twenty-five item resolutions
        and of defects D1 and D2.
  \item \texttt{Author Placeholder}. Drill protocol specification: score each
        examination item and cite the supporting lines in the source note.
        Internal working document, not distributed.
\end{enumerate}
\endgroup

% ===========================================================================
\appendix
% ===========================================================================

\section{Item Register: Twenty-Five Drill Items}
\label{app:register}

All twenty-five items supplied by the user were resolved against Module~7. For
each item the register records the accepted selection(s) and the rejected
option(s) with the note-grounded reason. ``\textsc{Resolved}'' means the note
determines the answer; the single ``\textsc{Ambiguous}'' entry corresponds to
open finding O3.

\small
\begin{longtable}{@{}p{0.5in}p{2.55in}p{3.3in}@{}}
\toprule
\textbf{\#} & \textbf{Accepted} & \textbf{Rejected, with reason} \\
\midrule
\endfirsthead
\toprule
\textbf{\#} & \textbf{Accepted} & \textbf{Rejected, with reason} \\
\midrule
\endhead
\bottomrule
\endfoot

1 & Salting, as the mechanism guaranteeing distinct digests for identical
    plaintext; a per-user salt guarantees divergence. &
    Symmetric encryption and asymmetric encryption: neither alters the digest,
    so neither can cause divergence. Peppering leaves the hashing function
    untouched and so leaves the digest unaltered. Resolved. \\

2 & User password manager (vault-file storage protected by a master password
    and a secret key file). &
    Enterprise vaulting: the described control is single-principal, so it is a
    user password manager. Resolved. \\

3 & Failure to enforce the possession element (a gate officer who admits on
    recognition). &
    Any option crediting the gate officer with enforcing possession:
    recognition is not possession. Resolved. \\

4 & Stolen digest files; reverse engineering (high-outcome). &
    Online brute force: classified low-outcome by the note. Resolved. \\

5 & Recommended against hardware password keys. &
    Hardware password keys: exposed to malware, which defeats the token's
    resistance to typed-input attacks. Resolved. \\

6 & Asymmetric cryptography, as the signing basis for assertion-based
    federation. &
    Symmetric alternatives: a relying party cannot verify a signature made
    with a shared secret it does not exclusively hold. Resolved. \\

7 & Key stretching \emph{and} Argon2, jointly, as digest protection. &
    None recorded. Resolved. \\

8 & Keystroke dynamics: requires neither specialized hardware nor additional
    user steps. &
    Any option attributing a hardware requirement or added user steps.
    Resolved. \\

9 & The time-based versus event-based one-time-password distinction is drawn
    from the note; the single selection the note supports is the one recorded,
    and the note was not pressed for a second. &
    Any second selection demanded by the item. The item demands two answers
    but the note supports only one. \textbf{Ambiguous (open finding O3).} \\

10 & Password reuse prohibition, as the sole defence against credential
     stuffing. &
    Any other sole defence: credential stuffing is defined by reuse, so
    permitting reuse leaves the characteristic intact. Resolved. \\

11 & Identity proofing and authentication, as the dual referents of the
     ``identity'' component. &
    Single-referent options treating the component as only one of the two.
     Resolved. \\

12 & Password length; reuse prohibition; minimum age (sound practice). &
    Expiration and an eight-character complexity floor. Resolved. \\

13 & Voice recognition; the relative difficulty of imitating cognitive
     biometrics. &
    Categorical assignment of the picture-password scheme; absolute
    spoof-resistance claims (the note supports a comparison, not a
    guarantee). Resolved. \\

14 & Genetic scanning maps to the \emph{biological} element. &
    The outwardly exhibited element: a genome is not a behaviour performed in
    front of a sensor. Resolved. \\

15 & Offline brute force and dictionary attacks are related by subsumption (a
     dictionary attack is a restricted case). &
    The inverted online-offline relation; the malformed comparison option.
     Resolved. \\

16 & Dictionary attack ranks as the highest-yield early tool in the
     eight-stage escalation sequence, ahead of hybrid, mask, and terminal
     brute force. &
    Ranking hybrid, mask, or terminal brute force above dictionary attack.
     Resolved. \\

17 & Knowledge-based possession, as the requirement for cognitive biometrics. &
    Any option treating the requirement as biometric-similarity rather than
    possession of the answer. Resolved. \\

18 & Account-lockout avoidance and alarm avoidance, as the motivations for
     password spraying; low-outcome classification enforced. &
    Reclassifying password spraying as high-outcome. Resolved. \\

19 & Dynamic role assignment \textrightarrow{} rule-based role assignment;
     attribute-combining flexibility \textrightarrow{} attribute-based control. &
    User-modifiability and conditional-structure claims. Resolved. \\

20 & Predictable substitution-and-append structure, yielding weakness. &
    Options reading the structure as unpredictability. Resolved. \\

21 & Passkeys: hardware storage and multi-factor combination. &
    Any option crediting either property to a different mechanism. Resolved. \\

22 & Statistical mask construction, as a rule attack. &
    Generic offline enumeration: the construction is statistical, hence a rule
    attack. Resolved. \\

23 & Parallel graphics processing, as the desirable brute-force platform. &
    Any option scoring a non-parallel platform as desirable for a massively
    parallel workload. Resolved. \\

24 & Security keys, over short-message authentication, on the ground of the
     absence of transmitted codes. &
    Short-message authentication: there is a transmitted code to intercept or
    relay. Resolved. \\

25 & Access control lists, matched to inefficient per-access permission
     checking. &
    Efficiency claims: the note documents access control list limitations as
    inefficient and enterprise-ungainly. Resolved. \\

\end{longtable}
\normalsize

\end{document}
